% META: {"id": "TEXP-ARI-EX-046", "chapitre": "TEXP-ARITHMETIQUE", "type_objet": "exercice", "capacites": ["TEXP-ARI-C1"], "capacites_codes": ["C1"], "methodes": ["M3"], "parcours": 2, "competences": ["raisonner","raisonner"], "duree_min": 18, "mode_creation": "ex_nihilo", "sources_inspiration": [], "fichier_tex": "chapitres/TEXP-ARITHMETIQUE/exercices/TEXP-ARI-EX-046.tex", "corrige_tex": "chapitres/TEXP-ARITHMETIQUE/corriges/TEXP-ARI-CO-046.tex", "status": "approved"}
% BEGIN-VERIFY
% from sympy import *
% x = symbols('x')
% a = symbols('a')
% f = symbols('f', cls=Function)
% # Theoreme: si f'' >= 0 sur I, alors f(x) >= f(a) + f'(a)(x-a)
% # Demonstration: g(x) = f(x) - f(a) - f'(a)(x-a)
% # g(a) = 0, g'(a) = 0, g''(x) = f''(x) >= 0 => g convexe => g >= 0
% # Verification sur f(x) = x^2, a = 2
% fx = x**2
% fa = 4
% fpa = 4
% tangente = fa + fpa*(x - 2)
% g = fx - tangente
% assert expand(g - (x - 2)**2) == 0
% assert g.subs(x, 2) == 0
% assert g.subs(x, 0) == 4
% assert g.subs(x, 5) == 9
% END-VERIFY
\begin{exercice}{TEXP-ARI-EX-046}{2}{18}
On veut demontrer le theoreme : \emph{Si $f'' \geq 0$ sur un intervalle $I$, alors pour tout $a \in I$ et tout $x \in I$ : $f(x) \geq f(a) + f'(a)(x - a)$.}
\begin{enumerate}
  \item Poser $g(x) = f(x) - f(a) - f'(a)(x-a)$. Que vaut $g(a)$ ? $g'(a)$ ?
  \item Exprimer $g''(x)$ et en deduire le signe de $g$.
  \item Conclure.
\end{enumerate}

\end{exercice}
